\documentclass[11pt,a4paper,UTF8]{ctexart}
\usepackage{geometry}
\geometry{left=18mm,right=18mm,top=24mm,bottom=24mm,headheight=22pt,footskip=12mm}
\usepackage{amsmath,amssymb,mathtools,bm,esint,extarrows,centernot}
\usepackage{xcolor}
\usepackage{tcolorbox}
\tcbuselibrary{skins,breakable}
\usepackage{fancyhdr}
\usepackage{titlesec}
\usepackage{hyperref}
\usepackage{tikz}
\usepackage{booktabs}
\usepackage{tabularx}

% --- Color Palette (Purple & DeepPink Academic Theme) ---
\definecolor{DeepPurple}{HTML}{4C1D95}
\definecolor{TitlePurple}{HTML}{581C87}
\definecolor{SubPurple}{HTML}{6B21A8}
\definecolor{DeepPink}{HTML}{FF1493}
\definecolor{LightPinkBg}{HTML}{FFF5F8}
\definecolor{LightPurpleBg}{HTML}{F8F5FF}
\definecolor{GoldAccent}{HTML}{D97706}
\definecolor{DarkText}{HTML}{1F2937}

\hypersetup{
    colorlinks=true,
    linkcolor=TitlePurple,
    citecolor=DeepPink,
    urlcolor=SubPurple,
    pdfborder={0 0 0}
}

% --- Typography & Spacing ---
\linespread{1.3}
\setlength{\parskip}{0.35em plus 0.1em minus 0.1em}
\setlength{\parindent}{0pt}

% --- Section Titles Styling ---
\titleformat{\section}
  {\Large\bfseries\color{TitlePurple}}{\thesection}{1em}{}
  [\vspace{1mm}{\color{TitlePurple!30}\hrule height 1pt}]
\titleformat{\subsection}
  {\large\bfseries\color{SubPurple}}{\thesubsection}{0.8em}{}
\titleformat{\subsubsection}
  {\normalsize\bfseries\color{DeepPurple}}{\thesubsubsection}{0.6em}{}

% --- tcolorbox Environments ---
% 1. Theorem / Formula / Method / Analysis Box (Deep Purple)
\newtcolorbox{thmbox}[1][]{
    enhanced,
    breakable,
    colback=LightPurpleBg,
    colframe=TitlePurple,
    arc=3pt,
    boxrule=0pt,
    leftrule=3.5pt,
    left=10pt,
    right=10pt,
    top=8pt,
    bottom=8pt,
    title=#1,
    coltitle=white,
    colbacktitle=TitlePurple,
    fonttitle=\bfseries\small,
    attach boxed title to top left={yshift=-2mm, xshift=4mm},
    boxed title style={boxrule=0pt, arc=2pt}
}

% 2. Solution / Formula / Derivation Box (DeepPink)
\newtcolorbox{solbox}[1][]{
    enhanced,
    breakable,
    colback=LightPinkBg,
    colframe=DeepPink,
    arc=3pt,
    boxrule=0pt,
    leftrule=3.5pt,
    left=10pt,
    right=10pt,
    top=8pt,
    bottom=8pt,
    title=#1,
    coltitle=white,
    colbacktitle=DeepPink,
    fonttitle=\bfseries\small,
    attach boxed title to top left={yshift=-2mm, xshift=4mm},
    boxed title style={boxrule=0pt, arc=2pt}
}

% 3. Learning Goals / Summary Box
\newtcolorbox{targetbox}[1][]{
    enhanced,
    breakable,
    colback=LightPurpleBg,
    colframe=DeepPurple,
    arc=3pt,
    boxrule=1pt,
    left=10pt,
    right=10pt,
    top=8pt,
    bottom=8pt,
    title=#1,
    coltitle=white,
    colbacktitle=DeepPurple,
    fonttitle=\bfseries\small,
    attach boxed title to top left={yshift=-2mm, xshift=4mm},
    boxed title style={boxrule=0pt, arc=2pt}
}

\begin{document}

% ------------------- 讲义封面 (Cover Page) -------------------
\begin{titlepage}
\thispagestyle{empty}
\begin{tikzpicture}[remember picture,overlay]
    \fill[DeepPurple] (current page.north west) rectangle ([yshift=-7cm]current page.north east);
    \draw[DeepPink, line width=3.5pt] ([yshift=-7cm]current page.north west) -- ([yshift=-7cm]current page.north east);
    \draw[GoldAccent, line width=1pt] ([yshift=-7.12cm]current page.north west) -- ([yshift=-7.12cm]current page.north east);
    
    \node[anchor=north east, opacity=0.12, text=white] at ([xshift=-1.5cm, yshift=-0.5cm]current page.north east) {
        \fontsize{70}{70}\selectfont\bfseries 2027
    };
\end{tikzpicture}

\vspace*{0.8cm}
\begin{center}
    {\color{white}\fontsize{26}{32}\selectfont\bfseries 2027 考研数学(二) 核心讲义}\\[0.6cm]
    {\color{white}\fontsize{13}{18}\selectfont\bfseries 全国硕士研究生招生考试 · 统考数学(二) 核心精讲全书}\\[1.5cm]
\end{center}

\vspace{3.5cm}
\begin{center}
    {\color{GoldAccent}\large\bfseries $\blacktriangleright$ 基础强化 · 核心题解 · 考点深水区 $\blacktriangleleft$}\\[0.8cm]
    {\color{TitlePurple}\fontsize{24}{28}\selectfont\bfseries [强化]一元函数积分学的概念与性质}\\[0.6cm]
    {\color{DeepPink}\large\bfseries —— 考研高阶冲刺与深水区专攻 ——}\\[1.5cm]
\end{center}

\vfill

\begin{center}
\begin{tcolorbox}[
    colback=white,
    colframe=DeepPurple,
    width=0.88\textwidth,
    arc=4pt,
    boxrule=1.2pt,
    halign=center,
    drop shadow=gray!30
]
    \vspace{3mm}
    {\bfseries\large 编著团队：EscoffierZhou}\\[2.5mm]
    {\bfseries\color{TitlePurple} 目标院校：中国科学院大学 (UCAS) 计算机专硕 (22408)}\\[2.5mm]
    {\small\color{gray} 备考铁律：手算真题数字 · 错题白纸重做 · 408 客观题为王}\\[2.5mm]
    {\small 编制版本：2027 届首发版 · \today}
    \vspace{2mm}
\end{tcolorbox}
\end{center}
\vspace{0.8cm}
\end{titlepage}

% ------------------- 目录与页面主体配置 -------------------
\pagestyle{fancy}
\fancyhf{}
\fancyhead[L]{\small\color{gray}2027 考研数学(二) · 核心讲义系列}
\fancyhead[R]{\small\color{TitlePurple}\nouppercase{\leftmark}}
\fancyfoot[C]{\small\color{gray}— 第 \thepage\ 页 —}
\fancyfoot[R]{\small\color{gray}UCAS 22408}
\renewcommand{\headrulewidth}{0.5pt}

\pagenumbering{Roman}
\tableofcontents
\newpage
\pagenumbering{arabic}


\begin{targetbox}[本章复习目标与核心知识图谱]
\begin{itemize}
  \item 目的[1]:攻克变限积分内含参变量的“拆分-换元-求导”深水区技巧
\end{itemize}
\end{targetbox}

\begin{targetbox}[本章复习目标与核心知识图谱]
\begin{itemize}
  \item 目的[2]:建立原函数存在性与定积分可积性的四象限逻辑矩阵与反例图谱
\end{itemize}
\end{targetbox}

\begin{targetbox}[本章复习目标与核心知识图谱]
\begin{itemize}
  \item 目的[3]:彻底搞清积分第一中值定理中开区间与闭区间的严格数学判据
\end{itemize}
\end{targetbox}

\begin{targetbox}[本章复习目标与核心知识图谱]
\begin{itemize}
  \item 目的[4]:掌握振荡反常积分的狄利克雷与阿贝尔审敛法，理清绝对收敛与条件收敛
\end{itemize}
\end{targetbox}

\begin{targetbox}[本章复习目标与核心知识图谱]
\begin{itemize}
  \item 目的[5]:掌握含对数与三角因子的高阶反常积分“阶数压制”秒杀速查模型
\end{itemize}
\end{targetbox}


\section{1.变限积分内含参变量的“拆分-换元”求导全景图谱}


\begin{equation*}
F(x) = \int_0^x (x - t)f(t)\,\mathrm{d}t
\end{equation*}

\begin{equation*}
F(x) = x\int_0^x f(t)\,\mathrm{d}t - \int_0^x tf(t)\,\mathrm{d}t
\end{equation*}
\begin{thmbox}[模型[1]: 线性参变量可拆分离型]

\textbf{[1]} 结构特征: 被积函数中的参变量 $x$ 仅以线性多项式形式作为因式出现：


\textbf{[2]} 标准处理SOP: 严禁直接代入求导！必须先将参变量 $x$ 线性拆开并提至积分号外，再使用乘积求导法则：


两端求导：


\begin{equation*}
F'(x) = \left[ 1 \cdot \int_0^x f(t)\,\mathrm{d}t + x f(x) \right] - x f(x) = \int_0^x f(t)\,\mathrm{d}t
\end{equation*}


二阶求导：


\begin{equation*}
F''(x) = f(x)
\end{equation*}


\end{thmbox}

\begin{equation*}
F(x) = \int_0^x f(x - t)\,\mathrm{d}t
\end{equation*}
\begin{thmbox}[模型[2]: 复合平移与比例不可拆型(必须换元前置)]

\textbf{[1]} 结构特征: 参变量 $x$ 与积分变量 $t$ 深度纠缠于被积函数复合内部，如 $f(x - t)$、$f(xt)$、$f(x^2 - t^2)$：


\textbf{[2]} 标准处理SOP: 严禁将 $x$ 视作常数强行求导！必须换元前置，把参变量 $x$ 转移到积分限上：

\textbf{[1]} 令 $u = x - t$，则 $t = x - u$，$\mathrm{d}t = -\mathrm{d}u$；

\textbf{[2]} 换限: 当 $t = 0$ 时 $u = x$；当 $t = x$ 时 $u = 0$；

\textbf{[3]} 整理被积表达式：


\begin{equation*}
F(x) = \int_x^0 f(u)(-\mathrm{d}u) = \int_0^x f(u)\,\mathrm{d}u
\end{equation*}


\textbf{[4]} 求导: 由微积分基本定理直接得 $F'(x) = f(x)$。

\end{thmbox}
\begin{thmbox}[模型[3]: 二次复合与双限变元型]

\textbf{[1]} 典型模型: $G(x) = \int_0^x t f(x^2 - t^2)\,\mathrm{d}t$。

\textbf{[2]} 换元手法: 凑微分 $\mathrm{d}(t^2)$，令 $u = x^2 - t^2$：

\textbf{[1]} 微分关系: $\mathrm{d}u = -2t\,\mathrm{d}t \implies t\,\mathrm{d}t = -\frac{1}{2}\mathrm{d}u$；

\textbf{[2]} 上下限: $t = 0 \to u = x^2$；$t = x \to u = 0$；

\textbf{[3]} 转化为标准定积分：


\begin{equation*}
G(x) = \int_{x^2}^0 f(u)\left(-\frac{1}{2}\mathrm{d}u\right) = \frac{1}{2}\int_0^{x^2} f(u)\,\mathrm{d}u
\end{equation*}


\textbf{[4]} 求导应用链式法则：


\begin{equation*}
G'(x) = \frac{1}{2} f(x^2) \cdot (x^2)' = x f(x^2)
\end{equation*}


\end{thmbox}

\section{2.原函数存在性与可积性的“四象限”逻辑矩阵}

\begin{thmbox}[原理[1]: 原函数存在与定积分可积的本质差异]

| \textbf{四象限分类} | \textbf{原函数存在性} | \textbf{定积分可积性} | \textbf{典型代表反例与核心机制} |
| :--- | :--- | :--- | :--- |
| \textbf{第一象限(双优)} | \textbf{存在} | \textbf{存在} | \textbf{闭区间连续函数} $f(x) \in C[a, b]$，连续必有原函数且必可积 |
| \textbf{第二象限(有原无积)} | \textbf{存在} | \textbf{不存在} | $f(x) = \begin{cases} 2x\sin\frac{1}{x^2} - \frac{2}{x}\cos\frac{1}{x^2}, & x \neq 0 \\ 0, & x = 0 \end{cases}$，原函数为 $x^2\sin\frac{1}{x^2}$ 处处可导，但在 $x=0$ 附近\textbf{无界}破坏可积性 |
| \textbf{第三象限(有积无原)} | \textbf{不存在} | \textbf{存在} | 含有\textbf{第一类跳跃间断点}的函数(如符号函数 $\text{sgn}(x)$)，有界可积；但由达布定理导函数不能跳跃，故\textbf{必无原函数} |
| \textbf{第四象限(双无)} | \textbf{不存在} | \textbf{不存在} | $f(x) = \frac{1}{x}$ 在 $[-1, 1]$ 上，含有无穷间断点必无原函数，无界必不可积 |

\end{thmbox}
\begin{thmbox}[原理[2]: 黎曼函数与狄利克雷函数的可积性深水区]

\textbf{[1]} 狄利克雷函数(Dirichlet): $D(x) = \begin{cases} 1, & x \in \mathbb{Q} \\ 0, & x \notin \mathbb{Q} \end{cases}$。

\textbf{[1]} 在任何区间上有界，但在任意小区间内上确界均为 1、下确界均为 0；

\textbf{[2]} 上下达布和恒有差距，在任何区间上均黎曼不可积；

\textbf{[3]} 但是 $D^2(x) \equiv 1$ 在任何闭区间上均可积(说明：可积函数的复合函数不一定可积)。

\textbf{[2]} 黎曼函数(Riemann): $R(x) = \begin{cases} \frac{1}{q}, & x = \frac{p}{q} \in (0, 1) \text{ 为既约真分数} \\ 0, & x \text{ 为无理数或 } 0, 1 \end{cases}$。

\textbf{[1]} 黎曼函数在无理点处连续，在有理点处全为可去间断点；

\textbf{[2]} 间断点集为可数集(零测度集)，因此黎曼函数在 $[0, 1]$ 上黎曼可积，且定积分值恒有：


\begin{equation*}
\int_0^1 R(x)\,\mathrm{d}x = 0
\end{equation*}


\end{thmbox}

\section{3.积分第一中值定理开闭区间的严格数学判定}

\begin{thmbox}[原理[1]: 介值定理法与拉格朗日法的证明差异]

积分第一中值定理表述为：设 $f(x)$ 在 $[a, b]$ 上连续，则存在中值点 $\xi$ 使得 $\int_a^b f(x)\,\mathrm{d}x = f(\xi)(b-a)$。

\textbf{[1]} 方法一(介值定理法):

由 $m(b-a) \leqslant \int_a^b f(x)\,\mathrm{d}x \leqslant M(b-a) \implies m \leqslant \frac{1}{b-a}\int_a^b f(x)\,\mathrm{d}x \leqslant M$。

由介值定理只能直接保证：
\begin{equation*}
\xi \in [a, b] \quad (\text{闭区间})
\end{equation*}


\textbf{[2]} 方法二(拉格朗日中值定理法):

构造变上限积分函数 $F(x) = \int_a^x f(t)\,\mathrm{d}t$。

因 $f(t)$ 连续，由微积分基本定理知 $F(x)$ 在 $[a, b]$ 上连续，在 $(a, b)$ 内可导，且 $F'(x) = f(x)$。

对 $F(x)$ 在 $[a, b]$ 上应用拉格朗日中值定理：


\begin{equation*}
F(b) - F(a) = F'(\xi)(b - a) \implies \int_a^b f(t)\,\mathrm{d}t = f(\xi)(b - a)
\end{equation*}


因拉格朗日中值定理结论严格落在开区间内部，故直接保证了：
\begin{equation*}
\xi \in (a, b) \quad (\text{严格开区间})
\end{equation*}


\end{thmbox}
\begin{thmbox}[定理[1]: 取得严格开区间的充要与充分条件]

\textbf{[1]} 充分条件: 只要 $f(x)$ 在 $[a, b]$ 上连续，必能找到开区间内部的点 $\xi \in (a, b)$ 满足积分中值公式。

\textbf{[1]} 若 $f(x)$ 是非常值连续函数: 最大值 $M$ 与最小值 $m$ 不可能在全区间恒相等，积分均值必定落在开区间 $(m, M)$ 内，由介值定理即可直接内插到开区间 $(a, b)$；

\textbf{[2]} 若 $f(x) \equiv C$ 为常数函数: 区间内任意一点(包括内部所有点)均显然满足 $f(\xi) = C$。

\end{thmbox}

\section{4.振荡反常积分的狄利克雷审敛法与条件收敛}

\begin{thmbox}[定理[1]: 狄利克雷判别法(Dirichlet's Test for Integrals)]

设函数 $f(x), g(x)$ 定义在 $[a, +\infty)$ 上：

\textbf{[1]} 变上限积分 $F(A) = \int_a^A f(x)\,\mathrm{d}x$ 在 $[a, +\infty)$ 上有界；

\textbf{[2]} $g(x)$ 在 $[a, +\infty)$ 上单调递减，且 $\lim\limits_{x \to +\infty} g(x) = 0$。

结论: 反常积分 $\int_a^{+\infty} f(x)g(x)\,\mathrm{d}x$ 必定收敛。

\end{thmbox}
\begin{thmbox}[定理[2]: 阿贝尔判别法(Abel's Test for Integrals)]

设函数 $f(x), g(x)$ 定义在 $[a, +\infty)$ 上：

\textbf{[1]} $\int_a^{+\infty} f(x)\,\mathrm{d}x$ 收敛；

\textbf{[2]} $g(x)$ 在 $[a, +\infty)$ 上单调有界。

结论: 反常积分 $\int_a^{+\infty} f(x)g(x)\,\mathrm{d}x$ 必收敛。

\end{thmbox}
\begin{thmbox}[模型[1]: 振荡基准积分 $\int_1^{+\infty} \frac{\sin x}{x^p}\,\mathrm{d}x$ 的全域收敛图谱]

\textbf{[1]} $p > 1$ 时:

由 $\left\vert{}\frac{\sin x}{x^p}\right\vert{} \leqslant \frac{1}{x^p}$，且 $\int_1^{+\infty} \frac{1}{x^p}\,\mathrm{d}x$ 收敛，故原积分绝对收敛。

\textbf{[2]} $0 < p \leqslant 1$ 时:

$\int_1^A \sin x\,\mathrm{d}x = \cos 1 - \cos A$ 在全区间上有界($\leqslant 2$)，而 $g(x) = \frac{1}{x^p}$ 单调趋于 0。由狄利克雷判别法知积分收敛。

但由三角倍角降幂公式 $\vert{}\sin x\vert{} \geqslant \sin^2 x = \frac{1-\cos 2x}{2}$ 可证其绝对值积分发散。故此时原积分条件收敛。

\textbf{[3]} $p \leqslant 0$ 时:

被积函数在 $x \to +\infty$ 时不趋向于 0，反常积分发散。

\end{thmbox}

\begin{equation*}
\text{P.V.} \int_{-\infty}^{+\infty} f(x)\,\mathrm{d}x = \lim_{A \to +\infty} \int_{-A}^A f(x)\,\mathrm{d}x
\end{equation*}
\begin{thmbox}[原理[1]: 柯西主值(P.V.)与反常积分收敛的不可互推性]

\textbf{[1]} 柯西主值定义:


\textbf{[2]} 本质陷阱:

考察反例 $f(x) = x$：


\begin{equation*}
\lim_{A \to +\infty} \int_{-A}^A x\,\mathrm{d}x = \lim_{A \to +\infty} 0 = 0
\end{equation*}


其对称主值存在且等于 0。

但常义反常积分 $\int_{-\infty}^{+\infty} x\,\mathrm{d}x = \int_{-\infty}^0 x\,\mathrm{d}x + \int_0^{+\infty} x\,\mathrm{d}x$ 两侧均发散，故原反常积分严格发散！

考研警示: 绝不能用对称区间的奇函数直接将两端发散的反常积分判定为收敛于 0。

\end{thmbox}

\section{5.含对数因子反常积分的“阶数压制”秒杀速查表}

对数函数 $\ln x$ 是“最慢增长”与“最弱发散”函数，在多项式幂函数面前，任何正微小量 $\varepsilon > 0$ 均可对其形成阶数压制。

\begin{thmbox}[速查表[1]: 无穷区间型 $\int_a^{+\infty} \frac{\ln^q x}{x^p}\,\mathrm{d}x \quad (a > 1, q \in \mathbb{R})$]

\textbf{[1]} 当 $p > 1$ 时: 对任意实数 $q$，积分必收敛(幂次 $p$ 占决定性主导地位)；

\textbf{[2]} 当 $p < 1$ 时: 对任意实数 $q$，积分必发散；

\textbf{[3]} 当 $p = 1$ 时: 转化为 $\int u^q\,\mathrm{d}u$，当且仅当 $q < -1$ 时收敛。

\end{thmbox}
\begin{thmbox}[速查表[2]: 瑕积分型 $\int_0^a x^p \vert{}\ln x\vert{}^q\,\mathrm{d}x \quad (0 < a < 1, q \in \mathbb{R})$]

\textbf{[1]} 当 $p > -1$ 时: 对任意实数 $q$，积分必收敛；

\textbf{[2]} 当 $p < -1$ 时: 对任意实数 $q$，积分必发散；

\textbf{[3]} 当 $p = -1$ 时: 写为 $\int_0^a \frac{\vert{}\ln x\vert{}^q}{x}\,\mathrm{d}x$，当且仅当 $q < -1$ 时收敛。
\end{thmbox}

\end{document}
